% !TEX root = ../main.tex
% File: chapters_part1/chap7_2.tex
% Nội dung cho Chương 7, Phần 2

\section{Các benchmark kinh điển và hiện đại}
\label{sec:benchmarks}

Một metric riêng lẻ chỉ đo một khía cạnh của mô hình. \textbf{Benchmark} là những bộ sưu tập bộ dữ liệu và tác vụ đã được tiêu chuẩn hóa, cho phép so sánh các mô hình trong cùng điều kiện và theo dõi tiến bộ của cả ngành.

\subsection{Kỷ nguyên Pre-LLM: GLUE và SuperGLUE}
\label{ssec:pre_llm_benchmarks}
Trước khi các LLM lớn với khả năng zero-shot ra đời, các mô hình (như BERT) thường được đánh giá bằng cách tinh chỉnh (fine-tune) trên một loạt tác vụ hiểu ngôn ngữ tự nhiên (NLU). GLUE và SuperGLUE là hai benchmark tiêu biểu của kỷ nguyên này.

\subsubsection{GLUE: bài kiểm tra năng lực ngôn ngữ tổng quát}
\begin{itemize}
    \item \textbf{Tên đầy đủ:} General Language Understanding Evaluation.
    \item \textbf{Giả định:} Một mô hình thực sự hiểu ngôn ngữ phải hoạt động tốt trên nhiều loại tác vụ khác nhau: suy luận, phân tích cảm xúc, nhận dạng sự tương đồng.
    \item \textbf{Cấu trúc:} GLUE gồm 9 tác vụ NLU, chia thành ba nhóm:
        \begin{itemize}
            \item suy luận ngôn ngữ tự nhiên (NLI): MNLI, QNLI, RTE;
            \item tương đồng ngữ nghĩa và diễn giải: MRPC, QQP, STS-B;
            \item phân loại câu đơn: SST-2 (cảm xúc), CoLA (tính đúng ngữ pháp).
        \end{itemize}
    \item \textbf{Đánh giá:} Mô hình được fine-tune riêng cho từng tác vụ; điểm cuối cùng là trung bình trên tất cả các tác vụ.
    \item \textbf{Tác động:} GLUE trở thành benchmark chính để đánh giá các mô hình như BERT, RoBERTa. Nhưng các mô hình nhanh chóng đạt tới và vượt qua hiệu năng của con người trên benchmark này vào khoảng năm 2019, nghĩa là nó đã trở nên quá dễ.
\end{itemize}

\subsubsection{SuperGLUE: nâng cấp độ khó}
\begin{itemize}
    \item \textbf{Động lực:} GLUE đã bị chinh phục, nên SuperGLUE được tạo ra với các tác vụ khó hơn, đòi hỏi suy luận phức tạp hơn.
    \item \textbf{Cấu trúc:} Cũng là một bộ sưu tập tác vụ đa dạng, nhưng tập trung hơn vào hỏi--đáp yêu cầu suy luận (BoolQ, MultiRC), suy luận nhân quả (COPA) và giải quyết đồng tham chiếu (WSC).
    \item \textbf{Tác động:} SuperGLUE trụ được lâu hơn, nhưng rồi các mô hình lớn hơn cũng dần chinh phục nó. Sự bão hòa của cả hai benchmark là một giới hạn của khuôn đánh giá bằng fine-tuning, và là lý do xuất hiện các phương pháp đánh giá mới.
\end{itemize}

\subsection{Kỷ nguyên LLM: đánh giá các khả năng nổi trội (Emergent Abilities)}
\label{ssec:llm_benchmarks}
Với các LLM có khả năng few-shot, khuôn đánh giá thay đổi: thay vì fine-tuning, các benchmark mới đo kiến thức và khả năng suy luận của mô hình trong cài đặt zero-shot và few-shot.

\subsubsection{MMLU: bài kiểm tra kiến thức đa chuyên ngành}
\begin{itemize}
    \item \textbf{Tên đầy đủ:} Massive Multitask Language Understanding.
    \item \textbf{Giả định:} Một LLM mạnh phải có kiến thức trên nhiều lĩnh vực chuyên môn, tương tự như các kỳ thi chuẩn hóa của con người.
    \item \textbf{Cấu trúc:} MMLU \cite{hendrycks2021mmlu} gồm 15\,908 câu hỏi trắc nghiệm 4 lựa chọn lấy từ các nguồn thi thực tế, bao phủ 57 môn học, từ các môn cơ bản (toán tiểu học) đến các môn chuyên ngành cấp đại học và chuyên nghiệp (luật, y, kỹ thuật).
    \item \textbf{Đánh giá:} Ở chế độ few-shot, thường là 5-shot. Prompt gồm 5 câu hỏi và câu trả lời mẫu từ một môn học, sau đó là câu hỏi cần trả lời.
    \item \textbf{Tác động:} MMLU trở thành một trong những benchmark chính để đo kiến thức của LLM. Nhưng đến 2024--2025 các mô hình hàng đầu đã vượt 90\% và benchmark gần như bão hòa, cộng thêm lo ngại nhiễm bẩn dữ liệu (mục~\ref{ssec:data_contamination}). Các phiên bản khó hơn như MMLU-Pro (10 lựa chọn, nhiều câu cần suy luận hơn) và GPQA (câu hỏi cấp tiến sĩ được thiết kế để khó tra cứu trên Google) được dùng thay thế.
\end{itemize}

\subsubsection{BIG-bench: thăm dò giới hạn của LLM}
\begin{itemize}
    \item \textbf{Tên đầy đủ:} Beyond the Imitation Game benchmark.
    \item \textbf{Giả định:} Cần một bộ sưu tập tác vụ đủ đa dạng để thăm dò các khả năng và những điểm yếu bất ngờ của LLM, điều mà các benchmark hiện có chưa làm được.
    \item \textbf{Cấu trúc:} 450 tác giả từ 132 tổ chức \cite{srivastava2023bigbench} cùng xây dựng 204 tác vụ, bao gồm logic, lý luận toán học, hiểu biết vật lý trực quan, nhận biết định kiến xã hội, và nhiều tác vụ khác mà máy tính trước đây chưa từng làm tốt.
    \item \textbf{Đánh giá:} Ở cả chế độ zero-shot và few-shot, kèm điểm của các chuyên gia con người làm mốc so sánh.
    \item \textbf{Tác động:} BIG-bench cho phép nghiên cứu cách năng lực thay đổi theo quy mô, trong đó có các “khả năng nổi trội” (emergent abilities) -- những khả năng chỉ bứt phá khi mô hình đạt một quy mô đủ lớn, thay vì cải thiện dần đều -- và việc thiên kiến xã hội thường tăng theo quy mô trong các ngữ cảnh mơ hồ. Tập con 23 tác vụ khó nhất, \textbf{BIG-Bench Hard (BBH)}, trở thành một chuẩn đo phổ biến cho Chain-of-Thought.
\end{itemize}

\subsubsection{GSM8K và MATH: đánh giá suy luận toán học}
\label{sssec:gsm8k_math}
\begin{itemize}
	\item \textbf{GSM8K} \cite{cobbe2021gsm8k}: 8{,}5 nghìn bài toán đố cấp tiểu học (7{,}5 nghìn huấn luyện, 1 nghìn kiểm tra), mỗi bài cần 2--8 bước tính số học cơ bản, đáp án là một con số. Đơn giản về kiến thức nhưng đòi hỏi suy luận nhiều bước không sai sót, nên nó là benchmark kinh điển của Chain-of-Thought. Đến nay đã gần như bão hòa với mô hình lớn.
	\item \textbf{MATH} \cite{hendrycks2021math}: 12\,500 bài toán thi đấu cấp trung học (AMC, AIME\ldots), trải trên 7 môn (Tiền đại số, Đại số, Lý thuyết số, Tổ hợp--Xác suất, Hình học, Đại số trung cấp, Tiền giải tích) và 5 mức độ khó, kèm lời giải từng bước viết bằng \LaTeX. Khi ra đời, ngay cả các mô hình lớn nhất cũng chỉ đạt điểm rất thấp; tập con 500 bài \textbf{MATH-500} (mục~\ref{ssec:prm}) hiện là chuẩn đo chính cho các mô hình suy luận, bên cạnh các đề thi mới như AIME để tránh nhiễm bẩn.
\end{itemize}
\textbf{Chấm điểm:} So khớp đáp án cuối cùng sau khi chuẩn hóa (ví dụ \texttt{1/2} và \texttt{0.5}, hay các biểu thức toán tương đương) -- chi tiết nhỏ này có thể làm lệch kết quả vài điểm phần trăm giữa các framework đánh giá khác nhau.

\subsubsection{TruthfulQA: đo lường sự trung thực}
\label{sssec:truthfulqa}
\begin{itemize}
	\item \textbf{Ý tưởng \cite{lin2022truthfulqa}:} mô hình được huấn luyện để bắt chước văn bản của con người sẽ bắt chước cả những \textbf{ngộ nhận phổ biến} của con người (“imitative falsehoods”).
	\item \textbf{Cấu trúc:} 817 câu hỏi trên 38 chủ đề (sức khỏe, pháp luật, tài chính, chính trị\ldots), được thiết kế có chủ đích để nhiều người trả lời sai do niềm tin sai lầm, ví dụ về mê tín hay thuyết âm mưu.
	\item \textbf{Kết quả:} Mô hình tốt nhất lúc công bố chỉ trung thực ở 58\% câu hỏi so với 94\% của con người, và mô hình càng lớn thường càng kém trung thực, vì chúng học thuộc các ngộ nhận phổ biến tốt hơn. Đây là ví dụ điển hình của \emph{inverse scaling}: tăng quy mô đơn thuần không giải quyết được vấn đề trung thực, mà cần căn chỉnh (InstructGPT cải thiện đáng kể trên TruthfulQA, mục~\ref{sssec:instructgpt}).
\end{itemize}

\subsubsection{HumanEval: đánh giá năng lực sinh mã nguồn}
\begin{itemize}
    \item \textbf{Tác vụ:} Viết và hoàn thiện mã nguồn (code generation).
    \item \textbf{Cấu trúc:} 164 bài toán lập trình, mỗi bài được đưa ra dưới dạng một docstring (mô tả hàm) và một vài test case.
    \item \textbf{Đánh giá:} Mô hình sinh ra phần thân của hàm; mã nguồn sinh ra được thực thi và kiểm tra xem có vượt qua các test case hay không. Metric chính là \textbf{Pass@k}, xác suất một giải pháp đúng được tạo ra trong $k$ lần thử.
    \item \textbf{Tác động:} HumanEval trở thành chuẩn đo năng lực lập trình của các mô hình như Codex, AlphaCode và các LLM hiện đại khác.
\end{itemize}

\begin{center}
\bookimage[1.0\textwidth]{llm_benchmarks_timeline}{Trục thời gian ngang từ 2016 đến 2024, trên đó đánh dấu năm ra đời của các benchmark tiêu biểu, mỗi benchmark là một vạch dựng lên kèm tên. Nền chia làm ba vùng màu theo thế hệ: vùng “hiểu ngôn ngữ” gồm SQuAD 1.1 (2016), SQuAD 2.0 và GLUE (2018), SuperGLUE (2019); vùng “kiến thức và suy luận” gồm MMLU, GSM8K, MATH (2021), BIG-bench (2022), HELM (2023); vùng “tác tử, kỹ thuật phần mềm” gồm AgentBench, GAIA và SWE-bench (2024). Khoảng cách giữa các mốc cho thấy mỗi thế hệ benchmark chỉ giữ được vai trò trong vài năm trước khi phải thay bằng thế hệ khó hơn.}
\captionof{figure}{Sơ đồ thời gian thể hiện sự phát triển của các benchmark}
\label{fig:llm_benchmarks_timeline}
\end{center}

Việc chuyển từ các benchmark như GLUE sang MMLU/BIG-bench thay đổi cả quy mô lẫn cách hiểu và đo “trí thông minh” của mô hình ngôn ngữ.

\subsection{HELM: đánh giá toàn diện thay vì một con số}
\label{ssec:helm}
Mỗi benchmark trên chỉ đo \emph{một} khía cạnh (thường là độ chính xác) trên \emph{một} kiểu tác vụ. Khi so sánh hai mô hình, các bài báo lại thường chọn những benchmark khác nhau, khiến việc so sánh gần như bất khả thi. \textbf{HELM} (Holistic Evaluation of Language Models) \cite{liang2022helm} của Stanford đề xuất một khung đánh giá có hệ thống dựa trên ba nguyên tắc:
\begin{enumerate}
	\item \textbf{Phân loại không gian đánh giá trước, rồi mới chọn tập con:} Định nghĩa rõ các \emph{kịch bản} (scenario = tác vụ + miền + ngôn ngữ, ví dụ “hỏi đáp về y tế bằng tiếng Anh”) và các \emph{thước đo}; thừa nhận công khai những gì \emph{chưa} được đánh giá.
	\item \textbf{Đa thước đo:} Với mỗi kịch bản, đo 7 thước đo: độ chính xác, hiệu chỉnh (calibration: mô hình có biết khi nào nó không chắc hay không), độ bền vững (robustness, khi đầu vào bị biến đổi nhẹ), công bằng, thiên kiến, độc hại, và hiệu quả (chi phí, độ trễ). Mô hình chính xác nhất chưa chắc đã tốt nhất trên các thước đo còn lại.
	\item \textbf{Chuẩn hóa điều kiện:} Mọi mô hình được đánh giá trên cùng kịch bản, cùng cách prompt, cùng số ví dụ few-shot.
\end{enumerate}
Phiên bản đầu tiên đánh giá 30 mô hình trên 16 kịch bản cốt lõi (cộng thêm 26 kịch bản mục tiêu). Trước HELM, trung bình mỗi mô hình chỉ được đánh giá trên 17{,}9\% các kịch bản cốt lõi này; HELM nâng con số đó lên 96\%. Ba nguyên tắc trên, minh bạch, đa chiều và tái lập được, là nền tảng của các bảng xếp hạng hiện nay.

\begin{keyidea}{Bối cảnh benchmark năm 2026}
Các benchmark “kinh điển” của kỷ nguyên LLM đầu (MMLU, GSM8K, HumanEval) đã gần bão hòa. Cộng đồng chuyển sang: (1) benchmark khó hơn như GPQA, MMLU-Pro, AIME, Humanity's Last Exam; (2) benchmark liên tục làm mới, dùng câu hỏi ra đời sau thời điểm thu thập dữ liệu huấn luyện, để chống nhiễm bẩn; (3) đánh giá bằng sở thích con người như Chatbot Arena (mục~\ref{sec:llm_as_a_judge}); và (4) đánh giá tác tử trên các nhiệm vụ nhiều bước thực tế như SWE-bench và GAIA (mục~\ref{sec:agent_evaluation}).
\end{keyidea}